\chapter{Gestion de crise}\label{ch:fm:crisis-management}

Le 31 octobre 2011, une procédure de liquidation s'ouvrait contre MF Global, un important futures commission merchant (FCM), et ses clients découvrirent que l'argent qui
devait leur revenir avait disparu. Le régulateur constata ensuite que, durant la dernière semaine d'octobre, la société avait utilisé les fonds ségrégués
des clients pour soutenir ses propres opérations et celles de ses sociétés affiliées~; un tribunal ordonna la restitution de 1.212~milliard de dollars en 2013, et les clients furent
finalement remboursés en totalité. Dans les premiers jours, les sociétés qui compensaient par son intermédiaire devaient savoir lesquelles de leurs positions subsistaient encore, où se trouvait leur
marge et qui pouvait la mobiliser. Une crise laisse à une société quelques heures, et ce qu'elle peut faire dans ces heures-là se décide bien en amont.

\section{La première heure d'un incident technologique}

\begin{definition}[Plan de gestion de crise, cellule de crise]\label{def:fm:crisis-management:plan}
Un \emph{plan de gestion de crise}\index{plan de gestion de crise} (crisis management plan) définit, pour des types de crise déterminés (un système de trading hors de contrôle, une tension
sur la liquidité, une défaillance de contrepartie, une fermeture de marché), qui décide, ce qu'il peut faire, à qui il le fait savoir et dans quel ordre, avec les procédures d'exploitation (runbook) et les listes de
contacts pour s'en acquitter. La \emph{cellule de crise}\index{cellule de crise} (crisis committee) est le petit groupe que le plan désigne pour prendre ces décisions, avec le pouvoir de
suspendre le trading, réduire les positions, déplacer le collatéral et tirer des financements sans les validations habituelles.
\end{definition}

La première heure d'un incident technologique relève du livre~13 et du livre~15~: le kill switch et l'annulation de masse (mass cancel) arrêtent les dégâts, le responsable d'incident (incident commander)
pilote la réponse à partir de la procédure d'exploitation, la bascule (failover) transfère le trading vers le site de reprise après sinistre (livre~14) si le site principal est perdu. Ce que la \omterm{def:fm:crisis-management:plan}{cellule de crise}
ajoute la décision métier qui entoure la décision technique~: faut-il reprendre le trading, quelles positions aplatir, ce qu'il faut dire aux contreparties, aux clients et au
régulateur, et combien de trésorerie l'incident a déjà coûté. Le plan de continuité d'activité du livre~14 et les règles de résilience opérationnelle fixent les
objectifs de reprise~; le \omterm{def:fm:crisis-management:plan}{plan de gestion de crise} décide qui agit lorsque ces objectifs sont sur le point d'être manqués.

\section{Tensions sur la liquidité~: appels de marge qu'on ne peut honorer}

Une crise de liquidité est le test de résistance (stress test) du chapitre~14, qui se déroule en heures au lieu de jours~: les exigences augmentent chez chaque courtier et chaque chambre de compensation à la fois,
les sources de financement qui semblaient fiables sont retirées, et les actifs qu'on pourrait vendre un jour normal prennent plus de temps ou coûtent plus cher. L'\omterm{def:fm:treasury-and-funding:horizon}{horizon de survie}
se mesure en heures, et chaque action a un délai.

\begin{definition}[Exercice de liquidité]\label{def:fm:crisis-management:drill}
Un \emph{exercice de liquidité}\index{exercice de liquidité} (liquidity drill) répète une crise de liquidité sur les comptes, les conditions et les sources de financement réels de la société~: un
scénario préétabli est joué heure par heure, la \omterm{def:fm:crisis-management:plan}{cellule de crise} choisit les actions avec leurs délais et leurs coûts réels, et l'exercice consigne l'\omterm{def:fm:treasury-and-funding:horizon}{horizon de survie}
et les décisions, à comparer avec le plan.
\end{definition}

\section{Quand une contrepartie fait défaut}

\begin{definition}[Actifs bloqués]\label{def:fm:crisis-management:trapped}
Les \emph{actifs bloqués}\index{actifs bloqués} (trapped assets) sont la trésorerie ou les titres d'une société détenus chez une contrepartie qui a fait défaut ou a été placée
en procédure d'insolvabilité, que la société ne peut utiliser tant qu'un administrateur ou un liquidateur ne les lui restitue pas, si jamais~: le collatéral déposé, l'excédent de marge, les titres
réhypothéqués (livre~1), ou soldes non ségrégués comme l'exigeaient les règles.
\end{definition}

Lorsqu'un prime broker ou un courtier compensateur fait défaut, trois choses se produisent à la fois. Le collatéral qu'il détient est bloqué, en totalité ou en partie, pendant des mois ou
des années (le livre~1 a raconté l'administration du broker-dealer européen de Lehman Brothers). Les positions qu'il détenait doivent être transférées ou rétablies
ailleurs, et le nouveau courtier exige immédiatement une marge pour elles. Et les propres contreparties de la société, à l'annonce de la nouvelle, peuvent relever leurs conditions.
La ségrégation des actifs des clients (livre~3) est la protection~; MF Global a montré que sa valeur dépend du respect de cette obligation par le courtier.

\begin{dated}{2026-09}{Les faits publics des exemples du chapitre}\label{dat:fm:crisis-management:record}
\textbf{MF Global}~: la procédure de liquidation a débuté le 31 octobre 2011~; la plainte de la CFTC reprochait qu'au cours de la dernière semaine d'octobre 2011 la société
avait utilisé les fonds ségrégués des clients pour soutenir ses propres opérations et celles de ses sociétés affiliées~; une ordonnance de consentement du 8 novembre 2013 exigeait que MF Global Inc. verse
1.212~milliard de dollars de restitution aux clients et une amende de 100~millions de dollars, et le liquidateur a ensuite procédé à des répartitions finales assurant une
restitution intégrale. \textbf{Marché des gilts, 2022}~: le 28 septembre 2022, la Banque d'Angleterre a annoncé des achats temporaires d'obligations d'État britanniques à long terme
pour rétablir des conditions de marché ordonnées, après que son Comité de politique financière eut relevé les risques pour la stabilité financière liés au dysfonctionnement de ce marché~;
les adjudications se sont poursuivies jusqu'au 14 octobre.
\end{dated}

\section{Tutoriel~: soixante-douze heures}\label{tut:fm:crisis-management:tut}

\begin{tutorial}
\textbf{Objectif.} Lancer un \omterm{def:fm:crisis-management:drill}{exercice de liquidité} de 72 heures sur une société du type de celle du chapitre~14~: un choc de volatilité, un prime broker qui fait défaut, un collatéral bloqué~; trouver
l'\omterm{def:fm:treasury-and-funding:horizon}{horizon de survie} et le jeu d'actions le moins coûteux qui le prolonge au-delà de 72 heures. \textbf{État final~:} le graphique horaire de la trésorerie avec et sans les
actions (la \cref{fig:fm:crisis-management:cash}) et les étapes du plan (la \cref{fig:fm:crisis-management:steps}).
\begin{tutsteps}
\item \textbf{Les comptes.} Le prime broker A détient 44~millions de dollars pour une exigence de 40~millions de dollars~; le prime broker B 33~millions de dollars pour 30~millions de dollars~;
le courtier compensateur 22~millions de dollars pour 20~millions de dollars. La \omterm{def:fm:treasury-and-funding:cash}{trésorerie non grevée} s'élève à 30~millions de dollars.
\item \textbf{Le choc.} La \omterm{def:fm:treasury-and-funding:house}{marge maison} (house margin) augmente de 60~\% dans la première journée (\texttt{firm.treasury.house\_path}, sans blocage)~; la chambre de compensation relève
la marge initiale de 30~\% au premier cycle quotidien et de 50~\% au second.
\item \textbf{La défaillance.} Le prime broker B fait défaut à l'heure 6~: 60~\% des 33~millions de dollars qu'il détient sont bloqués, le reste revient à l'heure 48, et ses
positions sont transférées au prime broker A.
\item \textbf{Les actions.} Vendre des fonds monétaires (10~millions de dollars après 4 heures, coût de 0.02~million de dollars)~; mettre en repo le portefeuille obligataire (10~millions de dollars après 8
heures, 0.1~million de dollars)~; tirer sur la ligne de crédit confirmée (20~millions de dollars après 2 heures, 0.3~million de dollars, retirée à l'heure 12)~; réduire les positions de 10~\%, 20~\% ou 30~\%
(après 3 heures, à des coûts d'impact de 1, 3 et 6~millions de dollars).
\item \textbf{Le plan.} \texttt{firm.crisisdrill.greedy} ajoute à chaque fois l'action dont le coût par heure gagnée est le plus faible, jusqu'à ce que l'horizon atteigne 72 heures.
\end{tutsteps}
\end{tutorial}

\omcode{code/firm/crisisdrill/firm_crisisdrill.py}{52}{90}{Le moteur de liquidité horaire~: appels réglés sur la trésorerie à mesure que les exigences augmentent, les positions du courtier défaillant
sont transférées, son collatéral est bloqué, et les actions s'achèvent après leurs délais.}\label{lst:fm:crisis-management:run}

\begin{omfigure}
\begin{tikzpicture}
\begin{axis}[width=10.5cm,height=5.4cm,xlabel={\small heure},ylabel={\small trésorerie non grevée (millions de \$)},xmin=0,xmax=71,ymin=-50,ymax=75,
  tick label style={font=\footnotesize},grid=major,grid style={gray!20},table/col sep=comma,
  legend cell align=left,legend style={font=\scriptsize,at={(0.5,-0.3)},anchor=north,/tikz/every even column/.append style={column sep=8pt}},legend columns=3]
\addplot[omDef,very thick,const plot] table[x=hour,y=baseline]{figdata/desk/27-crisis-management/cash.csv};
\addplot[omThm,thick,const plot] table[x=hour,y=plan]{figdata/desk/27-crisis-management/cash.csv};
\addplot[gray,dashed,const plot] table[x=hour,y=nofail]{figdata/desk/27-crisis-management/cash.csv};
\draw[black,thin] (axis cs:0,0) -- (axis cs:71,0);
\draw[gray,dotted] (axis cs:6,-50) -- (axis cs:6,75) node[pos=0.05,right,font=\scriptsize] {défaillance du courtier};
\legend{sans action,plan le moins coûteux,sans défaillance}
\end{axis}
\end{tikzpicture}

\omcaption{Soixante-douze heures de \omterm{def:fm:treasury-and-funding:cash}{trésorerie non grevée} dans le cadre de l'exercice~: sans action, la société est à court de liquidités à l'heure 6, au moment où les positions du courtier défaillant
doivent de nouveau être margées auprès du courtier restant~; avec le plan le moins coûteux, elle conserve au moins 11.8~millions de dollars. Sans la défaillance, le choc seul épuise la
trésorerie à l'heure 23. Données~: \texttt{fm\_crisis}.}\label{fig:fm:crisis-management:cash}
\end{omfigure}

Sans action, la société est à court de liquidités à l'heure 6. La défaillance fait le plus de dégâts~: 19.8~millions de dollars de collatéral sont bloqués chez le courtier défaillant, et ses
positions doivent être margées auprès du survivant dès maintenant~; le choc seul aurait épuisé la trésorerie à l'heure 23. Sur les 72 heures, les appels de marge totalisent 78.8
millions de dollars. Le plan glouton vend d'abord les fonds monétaires (l'horizon passe à l'heure 9 pour 20\,000 dollars), puis met les obligations en repo (heure 15), puis tire sur la
ligne de crédit avant son retrait (heure 24), et ce n'est qu'ensuite qu'on réduit de 10~\% les positions, ce qui permet à la société de dépasser 72 heures~: 1.42~million de dollars au total, dont l'essentiel
est l'impact de marché de la réduction. Les trois sources de trésorerie n'atteignent ensemble que l'heure 24~; sans la réduction, le second cycle quotidien de compensation épuiserait encore la
trésorerie.

\begin{omfigure}
\begin{tikzpicture}
\begin{axis}[width=10.5cm,height=5cm,xbar,bar width=9pt,xmin=0,xmax=80,xlabel={\small horizon de survie (heures)},ytick={0,1,2,3,4},
  yticklabels={sans action,{+ vente de fonds monétaires},{+ repo des obligations},{+ tirage de la ligne},{+ réduction de 10~\% des positions}},y dir=reverse,
  ymin=-0.6,ymax=4.6,tick label style={font=\footnotesize},grid=major,grid style={gray!20},table/col sep=comma,nodes near coords,
  nodes near coords style={font=\scriptsize}]
\addplot[fill=omDef!55,draw=omDef] table[y=pos,x=horizon]{figdata/desk/27-crisis-management/steps.csv};
\draw[gray,dashed] (axis cs:72,-0.6) -- (axis cs:72,4.6);
\end{axis}
\end{tikzpicture}

\omcaption{L'\omterm{def:fm:treasury-and-funding:horizon}{horizon de survie} au fur et à mesure que le plan glouton ajoute des actions, les moins coûteuses par heure gagnée d'abord~; coûts cumulés de 0, 0.02, 0.12, 0.42 et 1.42 million de dollars.
Données~: \texttt{fm\_crisis.steps}.}\label{fig:fm:crisis-management:steps}
\end{omfigure}

\begin{proposition}[Davantage de ressources ne raccourcit jamais l'horizon]\label{prop:fm:crisis-management:monotone}
Dans le modèle de l'\omterm{def:fm:crisis-management:drill}{exercice de liquidité}, ajouter une action qui apporte de la trésorerie ou réduit les exigences, sans modifier aucune autre action, ne rend jamais
l'\omterm{def:fm:treasury-and-funding:horizon}{horizon de survie} plus court.
\end{proposition}

\begin{proof}
La trésorerie à chaque heure est la trésorerie initiale, plus la trésorerie des actions achevées, plus les restitutions d'excédent, moins les appels. Une action de trésorerie ajoute un montant non négatif
à partir de son achèvement et ne modifie aucun appel. Une réduction des positions abaisse chaque exigence ultérieure, si bien que chaque appel ultérieur n'est pas supérieur (les prime brokers retiennent
l'excédent, et le compte de compensation le restitue, ce qui ne fait qu'ajouter de la trésorerie). La trésorerie est donc au moins aussi élevée à chaque heure, et la première heure négative n'est pas
plus tôt.
\end{proof}

\begin{remark}[Le glouton n'est pas optimal]\label{rem:fm:crisis-management:greedy}
Choisir selon le coût par heure gagnée est une heuristique. Les heures ne sont pas additives~: une action à long délai ne rapporte rien seule et beaucoup après
que d'autres ont acheté le temps dont elle a besoin, et une ligne retirée à l'heure 12 doit être tirée tôt ou jamais. Pour un plan de quelques actions, une
recherche exhaustive est peu coûteuse~; la valeur de l'exercice tient à ce que la \omterm{def:fm:crisis-management:plan}{cellule de crise} s'exerce au choix, et non à l'algorithme.
\end{remark}

\section{Exercices de liquidité et procédures types}

\begin{method}[Conduire un exercice de liquidité]\label{met:fm:crisis-management:drill}
\begin{enumerate}
\item Rédiger le scénario à partir des comptes et des termes propres à la firme~: quel courtier fait défaut, quels actifs sont bloqués, quelles lignes sont retirées, comment les marges montent, heure par
heure.
\item Remettre à la \omterm{def:fm:crisis-management:plan}{cellule de crise} les véritables listes de contacts, les véritables autorités et les délais réels de chaque action, et lancer l'exercice sans préavis.
\item Consigner chaque décision avec son heure~; calculer l'\omterm{def:fm:treasury-and-funding:horizon}{horizon de survie} atteint, à comparer à celui du plan.
\item Corriger ce que l'exercice a mis au jour~: une ligne qui exigeait deux signatures que personne ne pouvait joindre, un virement qui a manqué l'heure limite, une procédure d'exploitation (runbook) qui supposait
un système qui était en panne.
\item Répéter au moins une fois par an, et après tout changement de courtier, de ligne de financement ou de stratégie.
\end{enumerate}
\end{method}

Les recommandations de 2024 du Conseil de stabilité financière sur la préparation à la liquidité face aux appels de marge et aux appels de collatéral, résumées au chapitre 14, demandent
aux sociétés non bancaires exactement ceci~: des tests de résistance aux appels de marge dans des scénarios extrêmes mais plausibles, et des plans pour y faire face. Un exercice est le lieu où le
plan se confronte à l'horloge.

\section{Construction~: l'exercice de crise}\label{bld:fm:crisis-management:crisisdrill}

\begin{build}
\bfield{Rôle} Un \omterm{def:fm:crisis-management:drill}{exercice de liquidité} horaire~: des scénarios avec chocs, une contrepartie défaillante, des \omterm{def:fm:crisis-management:trapped}{actifs bloqués} et des lignes retirées~; l'\omterm{def:fm:treasury-and-funding:horizon}{horizon de survie}~;
et un plan d'actions choisi selon le coût par heure gagnée.
\bfield{Interface} \texttt{firm.crisisdrill}~: \texttt{Account}, \texttt{Scenario}, \texttt{Action}, \texttt{run}, \texttt{greedy}~; la trajectoire de la \omterm{def:fm:treasury-and-funding:house}{marge maison}
issue de \texttt{firm.treasury}.
\bfield{Règles} Les appels sont payés à l'heure où ils surviennent~; un compte défaillant s'arrête et restitue son collatéral non bloqué à l'heure annoncée~; une action
n'est disponible que jusqu'à sa date limite.
\bfield{Tests d'acceptation} \texttt{code/firm/crisisdrill/tests/}~: appels et horizon sur des chiffres de trésorerie disponible~; défaillance, blocage et virement~; le choix glouton
et une ligne retirée.
\bfield{Pour aller plus loin} Recherche exhaustive sur les ensembles d'actions~; délais incertains~; le modèle de marge propre à la chambre de compensation, tiré de \texttt{firm.initmargin}~; plusieurs
défaillances.
\end{build}

\begin{omsources}
\item Communiqués de presse de la CFTC 6776-13 (2013), 6904-14 (2014) et 7508-17 (2017) sur MF Global.
\item Bank of England, «~Bank of England announces gilt market operation~», 28 septembre 2022.
\end{omsources}

\section{Exercices}

\begin{exercise}[$\star$]\label{exo:fm:crisis-management:1}
Quel montant de collatéral est bloqué chez le courtier défaillant, et combien revient à l'heure 48~?
\end{exercise}

\begin{exercise}[$\star$]\label{exo:fm:crisis-management:2}
Définir les \omterm{def:fm:crisis-management:trapped}{actifs bloqués} et donner trois façons dont les actifs d'une firme peuvent être bloqués chez un courtier défaillant.
\end{exercise}

\begin{exercise}[$\star$]\label{exo:fm:crisis-management:3}
À quelle heure la firme manquerait-elle de trésorerie sans la défaillance, et pourquoi si tard~?
\end{exercise}

\begin{exercise}[$\star\star$]\label{exo:fm:crisis-management:4}
Pourquoi le plan glouton ne tire-t-il la ligne engagée qu'en troisième position, et que se passerait-il si la ligne exigeait deux heures d'approbations auprès d'une personne qui est
injoignable jusqu'à l'heure 12~?
\end{exercise}

\begin{exercise}[$\star\star$]\label{exo:fm:crisis-management:5}
Démontrer la \cref{prop:fm:crisis-management:monotone}. Resterait-elle valable si la réduction des positions entraînait des pertes payées en espèces~?
\end{exercise}

\begin{exercise}[$\star\star$]\label{exo:fm:crisis-management:6}
Que devaient savoir les clients de MF Global dans les premiers jours, et quels registres le leur auraient-ils indiqué~?
\end{exercise}

\begin{exercise}[$\star\star\star$]\label{exo:fm:crisis-management:7}
\emph{Programmation.} Lancer l'exercice avec les trois actions de trésorerie et sans réduction. Quel est l'\omterm{def:fm:treasury-and-funding:horizon}{horizon de survie}, et quel est le niveau de trésorerie le plus bas~?
\end{exercise}

\begin{exercise}[$\star\star\star$]\label{exo:fm:crisis-management:8}
\emph{Trouver la faille.} «~Nous avons 30~millions de dollars de trésorerie et 40~millions de dollars de lignes~: nous pouvons survivre à tout appel de marge.~»
\end{exercise}

\section{Problème~: soixante-douze heures}

\begin{problem}[{Problème du week-end --- soixante-douze heures}]\label{pb:fm:crisis-management:1}
Un \omterm{def:fm:the-risk-management-function:cro}{directeur des risques} (CRO) doit conduire le premier \omterm{def:fm:crisis-management:drill}{exercice de liquidité} de la firme et présenter au conseil d'administration ce qu'il en ressort.

\medskip\noindent\textbf{Partie I --- Le plan.}
\begin{enumerate}
\item Définir un \omterm{def:fm:crisis-management:plan}{plan de gestion de crise} et une \omterm{def:fm:crisis-management:plan}{cellule de crise}.
\item Qu'est-ce que la cellule apporte à la réponse à un incident technologique~?
\item Définir un \omterm{def:fm:crisis-management:drill}{exercice de liquidité}.
\item Définir les \omterm{def:fm:crisis-management:trapped}{actifs bloqués}.
\end{enumerate}

\medskip\noindent\textbf{Partie II --- Les faits.}
\begin{enumerate}[resume]
\item Résumer les faits publics concernant MF Global.
\item Que fit la Banque d'Angleterre en septembre 2022, et pourquoi~?
\item Qu'est-ce qui protège les actifs des clients chez un courtier, et qu'a montré MF Global à ce sujet~?
\item Que demandent aux sociétés non bancaires les recommandations internationales de 2024~?
\end{enumerate}

\medskip\noindent\textbf{Partie III --- L'exercice.}
\begin{enumerate}[resume]
\item Décrire le scénario heure par heure.
\item Donner l'\omterm{def:fm:treasury-and-funding:horizon}{horizon de survie} sans action et sans la défaillance.
\item Donner les étapes, les horizons et les coûts du plan glouton.
\item Énoncer et démontrer la \cref{prop:fm:crisis-management:monotone}.
\item Pourquoi le plan glouton n'est-il pas nécessairement le moins coûteux~?
\end{enumerate}

\medskip\noindent\textbf{Partie IV --- Le rapport.}
\begin{enumerate}[resume]
\item Quelles faiblesses un exercice réel ferait-il probablement apparaître que le modèle ne montre pas~?
\item Que changeriez-vous dans les dispositions de la société avec ses courtiers~?
\item Quelle doit être la taille du \omterm{def:fm:treasury-and-funding:cash}{coussin de liquidité}~?
\item À quelle fréquence l'exercice doit-il avoir lieu, et qui ne doit pas être prévenu à l'avance~?
\item Que rapporteriez-vous au conseil d'administration~?
\item Énoncer le \emph{résultat nommé}~: l'\omterm{def:fm:treasury-and-funding:horizon}{horizon de survie} sous l'exercice, et le coût de l'ensemble d'actions le moins coûteux qui le prolonge jusqu'à 72~heures.
\item En deux phrases, rédiger la recommandation.
\end{enumerate}
\end{problem}

\section{Questions d'entretien}

\begin{interviewq}[$\star$ \iqroles{développeur}]\label{iq:fm:crisis-management:1}
Une stratégie envoie des ordres qu'elle ne devrait pas envoyer. Que faites-vous dans les cinq premières minutes~?
\end{interviewq}

\begin{interviewq}[$\star$ \iqroles{risque}]\label{iq:fm:crisis-management:2}
Votre prime broker vient d'annoncer qu'il est placé sous administration. Que vérifiez-vous en premier~?
\end{interviewq}

\begin{interviewq}[$\star\star$ \iqroles{risque, trader}]\label{iq:fm:crisis-management:3}
Les appels de marge dépassent votre trésorerie de 10~millions de dollars demain matin. Listez vos options par ordre de priorité.
\end{interviewq}

\begin{interviewq}[$\star\star$ \iqroles{développeur, risque}]\label{iq:fm:crisis-management:4}
Comment concevriez-vous un \omterm{def:fm:crisis-management:drill}{exercice de liquidité} qui teste les personnes autant que les modèles~?
\end{interviewq}

\begin{interviewq}[$\star\star$ \iqroles{trader}]\label{iq:fm:crisis-management:5}
Pourquoi réduire ses positions en période de crise peut-il coûter plus qu'il ne rapporte~?
\end{interviewq}

\begin{interviewq}[$\star\star\star$ \iqroles{chercheur, risque}]\label{iq:fm:crisis-management:6}
Formulez le choix des actions de crise comme un problème d'optimisation, et expliquez pourquoi une règle gloutonne peut échouer.
\end{interviewq}
